\documentclass[11pt,a4paper]{article}
\usepackage[T1]{fontenc}
\usepackage[utf8]{inputenc}
\usepackage{lmodern,amsmath,amssymb,amsthm,mathtools,booktabs,array,microtype}
\usepackage[margin=27mm]{geometry}
\usepackage{xcolor,listings,xurl}
\usepackage[colorlinks=true,linkcolor=blue!45!black,urlcolor=blue!45!black,citecolor=blue!45!black]{hyperref}
\usepackage{fancyhdr}
\pagestyle{fancy}\fancyhf{}
\fancyhead[L]{\small Higher-dimensional research checkpoint}
\fancyhead[R]{\small No original conjecture resolved}
\fancyfoot[C]{\thepage}\setlength{\headheight}{14pt}
\setlength{\parindent}{0pt}\setlength{\parskip}{5pt}
\newtheorem{theorem}{Theorem}
\newtheorem{lemma}[theorem]{Lemma}
\newtheorem{proposition}[theorem]{Proposition}
\newcommand{\F}{\mathbb F_5}
\newcommand{\D}{\mathbb D}
\newcommand{\T}{\mathbb T}
\newcommand{\supp}{\operatorname{supp}}
\newcommand{\diag}{\operatorname{diag}}
\newcommand{\conv}{\operatorname{conv}}
\hypersetup{pdftitle={Higher-Dimensional Research Checkpoint: Exact Barriers and Unresolved Targets},pdfauthor={Research checkpoint prepared for Keio},pdfsubject={Unresolved named conjectures and proved construction barriers}}
\lstset{language=Python,basicstyle=\ttfamily\scriptsize,columns=fullflexible,
keepspaces=true,breaklines=true,showstringspaces=false,
keywordstyle=\color{blue!50!black},commentstyle=\color{black!55},
frame=single,rulecolor=\color{black!20},numbers=left,
numberstyle=\tiny\color{black!50},numbersep=6pt,xleftmargin=12pt}
\title{Higher-Dimensional Research Checkpoint:\\
Exact Barriers and Unresolved Targets}
\author{Research checkpoint prepared for Keio}
\date{8 November 2026}

\begin{document}
\maketitle
\begin{abstract}
This report records attempts on the Marcus--de Oliveira determinantal
conjecture, DeVos's matrix choosability conjecture, and the Krzy\.{z}
coefficient conjecture. None is resolved here. The proved statements are
more limited: a translate-covering lemma over $\mathbb F_5$ shows why an
entire family of affine block constructions cannot produce a DeVos
counterexample; a gain-graph argument gives a second affirmative structured
class in arbitrary even dimension; and an explicit entropy calculation
invalidates a naive extension of a hypothesis-dependent inequality in the
Krzy\.{z} programme. Numerical and finite search records are explicitly
separated from proofs. No claim of historical priority, external peer
review, or resolution of a previously open named conjecture is made.
\end{abstract}

\section{Outcome and exact scope}

\textbf{Number of newly resolved original named targets in this report: zero.}
The purpose of this document is to preserve exact statements, valid
arguments, and identified gaps so that subsequent work does not repeat a
failed construction or misidentify an existing result as new.

The original goals require all dimensions or all coefficient indices.
Checking a finite list of instances, proving a restricted matrix class, or
excluding a proposed proof strategy does not satisfy those quantifiers.
The elementary results proved below are unconditional within their stated
scope, but their novelty has not been established by a complete literature
audit. They are not substitutes for the original conjectures.

\section{The Marcus--de Oliveira determinantal conjecture}

For normal complex $n\times n$ matrices $A,B$, with eigenvalues
$a_1,\ldots,a_n$ and $b_1,\ldots,b_n$, the original conjecture asks whether
\begin{equation}\label{eq:moc}
 \det(A+B)\in
 \conv\left\{\prod_{i=1}^n(a_i+b_{\sigma(i)}):\sigma\in S_n\right\}
\end{equation}
for every $n$. The problem originates in work of Marcus and de Oliveira;
the history and known classes are surveyed in \cite{MOCsurvey}.
The current partial manuscript \cite{Holden} explicitly retains the
unrestricted target.

After a common unitary change of basis, write $A=D_a$ and
$B=UD_bU^*$, where $D_a,D_b$ are diagonal and $U$ is unitary. Put
\[
 d_U(a,b)=\det(D_a+UD_bU^*),\qquad
 z_\sigma(a,b)=\prod_i(a_i+b_{\sigma(i)}).
\]
Principal-minor expansion and Cauchy--Binet give the established identity
\begin{equation}\label{eq:minor}
 d_U(a,b)=\sum_{|I|=|J|}a_{I^c}b_J\,|\det U[I,J]|^2.
\end{equation}
It is tempting to find nonnegative weights $w_\sigma$, independent of
$a,b$, satisfying
\[
 |\det U[I,J]|^2
 =\sum_\sigma w_\sigma\,\mathbf1_{\sigma(I)=J},
 \qquad \sum_\sigma w_\sigma=1.
\]
This would prove \eqref{eq:moc} for the given $U$ and all spectra.
However, this stronger common-weight property is already known to fail.
Its failure does not refute \eqref{eq:moc}, which permits the convex
weights to depend on the spectra. The signed-weight reductions in
\cite{Holden} are prior work, not results of this report.

An actual counterexample would require exact data $U,a,b,\eta$, with $U$
unitary and
\begin{equation}\label{eq:separation}
 \Re\bigl(\eta d_U(a,b)\bigr)
 >\max_{\sigma\in S_n}\Re\bigl(\eta z_\sigma(a,b)\bigr).
\end{equation}
The separating-functional criterion is equivalent to being outside the
finite convex hull. No data satisfying it were obtained.

\paragraph{Exploratory computation.}
An order-four search sampled 400 orthogonal or unitary matrices, selected
30 numerically failing the common-weight feasibility test, and ran 600
spectral optimizations. The objective frequently collapsed to degenerate
spectra. A second search sampled 1000 matrices plus a prior rational
quaternion example, selected 15 matrices from 48 positive-obstruction
candidates, and ran 225 optimizations with the variance of the real
permutation vertices fixed at one. The best recorded separation was
approximately $-5.5061\times10^{-5}$, not positive. Its spectrum had three
equal eigenvalues on one side, approaching an already affirmative class.
These searches prove neither a finite exclusion nor a global bound.

\section{DeVos's matrix choosability conjecture}

The author-posted formulation \cite{DeVos} says that, for a prime $p$,
every invertible matrix $A\in\mathbb F_p^{n\times n}$ is
$(k+2,p-k)$-choosable. Explicitly, for arbitrary sets
\[
 |X_i|=k+2,\qquad |Y_i|=p-k,
\]
there should be an $x\in\prod_i X_i$ with $Ax\in\prod_i Y_i$.
We consider $1\le k\le p-2$.

At $p=5,k=1$, every three-element set is affine-equivalent to
\[
 D=\{0,1,-1\}\subset\F.
\]
Indeed, an affine bijection can send its complementary pair to
$\{2,-2\}$. Coordinatewise affine changes preserve invertibility of the
linear part. Each four-element $Y_i$ excludes one value. Thus a
counterexample to this particular original case is equivalent to an
invertible $A$ and an offset $b$ for which the independent affine
hyperplanes
\[
 (Ax)_i=b_i\qquad(1\le i\le n)
\]
cover $D^n$. This reduction does not settle other primes or other $k$.

\subsection{A forced relation and its exact limitation}

All arithmetic in this section is in $\F$. Let
\begin{equation}\label{eq:gadget}
 A_0=\begin{pmatrix}2&0&2\\1&4&0\\3&3&2\end{pmatrix},
 \qquad L=(0,1,1).
\end{equation}
Then $\det A_0=-2=3$ and $LA_0^{-1}=(3,4,0)$.
Scaling the first and third rows shows that $A_0(a,b,c)^T$ has no zero
coordinate exactly when
\[
 a+c\ne0,\qquad a-b\ne0,\qquad a+b-c\ne0.
\]
For $(a,b,c)\in D^3$ these conditions hold exactly at
\begin{equation}\label{eq:six}
 (a,b,-b),\qquad a,b\in D,\quad a\ne b.
\end{equation}
To check the forcing step, if $b+c\ne0$, the three forbidden values
$-c,b,c-b$ contain every element of $D$. If $c=-b$, the only additional
forbidden value besides $b$ is $-2b$, which is outside $D$ when
$b\ne0$. Thus there are six survivors, all satisfying $Lx=0$.

One might try to connect copies through their forced $L$-values to create
an inconsistency. The next lemma explains why this construction fails.

\begin{lemma}[Translate covering]\label{lem:translate}
The set
\[
 K=\{z\in\F^3:Lz=0,\ A_0z\in(\F^\times)^3\}
   =\{(u,v,-v):u\ne v,\ u\ne3v\}
\]
meets every translate of $D^3$ in $\F^3$.
\end{lemma}
\begin{proof}
Consider a translate $h+D^3$. Choose
\[
 v\in(h_2+D)\cap(-h_3-D).
\]
The intersection is nonempty because each set has three elements in a
five-element group. Choose $u\in h_1+D$ avoiding the at most two values
$v,3v$. Then $(u,v,-v)\in K\cap(h+D^3)$.
\end{proof}

\begin{lemma}[Affine local realization]\label{lem:local}
For every $a\in\F^3$ there is an $x\in D^3$ such that
\[
 A_0x+a\in(\F^\times)^3,
 \qquad Lx=-LA_0^{-1}a.
\]
\end{lemma}
\begin{proof}
Apply Lemma~\ref{lem:translate} to $A_0^{-1}a+D^3$, and write the
resulting point as $A_0^{-1}a+x\in K$.
\end{proof}

\begin{theorem}[An arbitrary-dimensional affine block class]\label{thm:network}
Let $m\ge1$, let $\mathcal D=\diag(A_0,\ldots,A_0)$ have $m$ blocks,
and let $R:\F^{3m}\to\F^m$ apply $L$ to each block. For any matrix
$C\in\F^{3m\times m}$, suppose
\[
 A=\mathcal D+CR
\]
is invertible. For every $b\in\F^{3m}$ there exists $x\in D^{3m}$
such that $Ax+b$ has no zero coordinate.
\end{theorem}
\begin{proof}
The determinant identity
\[
 \det(\mathcal D+CR)=\det(\mathcal D)
       \det(I_m+R\mathcal D^{-1}C)
\]
shows that the reduced system
\begin{equation}\label{eq:reduced}
 (I_m+R\mathcal D^{-1}C)r=-R\mathcal D^{-1}b
\end{equation}
has a unique solution. Put $a=Cr+b$, divided into three-coordinate
blocks $a_i$. Apply Lemma~\ref{lem:local} independently to obtain
$x_i\in D^3$ with $A_0x_i+a_i$ nowhere zero and
$Lx_i=-LA_0^{-1}a_i$. Equation~\eqref{eq:reduced} says precisely that
these latter values are $r_i$. Hence $Rx=r$, so
$Ax+b=\mathcal D x+Cr+b$ is nowhere zero.
\end{proof}

This proves an affirmative structured class in every dimension $3m$.
It also excludes every proposed affine counterexample network whose
cross-block interaction uses only the observable $L$.
Arbitrary invertible matrices do not have to belong to this class.

\subsection{A second structured class from one-cycle graphs}

The following argument formalizes another unsuccessful counterexample
construction. Its ingredients are elementary list coloring and the Schur
complement; no priority claim is made for the resulting class.

\begin{lemma}[A cycle with nonzero additive gain]\label{lem:cycle}
Let each vertex of a cycle have a list of at least two elements of $\F$.
On each oriented edge impose the condition that the difference of the
endpoint colors must avoid one specified value. If the signed sum of
these forbidden values around the cycle is nonzero, a list coloring
satisfying every condition exists.
\end{lemma}
\begin{proof}
An edge forbids a partial matching between its two endpoint lists.
If some list has at least three colors, leave that vertex until last
and greedily color the remaining path. The last vertex must avoid at
most two values. If every list has two colors but some forbidden
matching is imperfect, start with an unmatched color at an endpoint
of that edge and greedily color the path obtained by removing the edge.
The removed edge cannot be violated.

Otherwise every forbidden matching is perfect between two-element
sets. Going once around the cycle composes additive translations,
and therefore translates the starting two-element set by the nonzero
total gain while preserving it. This is impossible: any nonzero
translation of $\F$ has a single orbit of length five. These cases
exhaust all possibilities.
\end{proof}

\begin{theorem}[A unicyclic incidence class]\label{thm:unicyclic}
Let $G$ be a finite connected simple graph with $m$ vertices and $m$
edges, and let $E\in\F^{m\times m}$ be an oriented edge--vertex
incidence matrix. Let $U,V\in\F^{m\times m}$ and suppose
\[
 A=\begin{pmatrix}I_m&U\\E&V\end{pmatrix}
\]
is invertible. For every $b\in\F^{2m}$ there are $x,z\in D^m$ such
that $A(x,z)^T-b$ has no zero coordinate.
\end{theorem}
\begin{proof}
Put $W=V-EU$, which is invertible by the Schur complement.
The unique cycle of $G$ gives a nonzero signed cycle vector $c$ with
$cE=0$. Thus $cW\ne0$. Write $b=(b_{\rm top},b_{\rm bottom})$.
Choose $z\in D^m$ so that
\[
 c\bigl(Wz+Eb_{\rm top}-b_{\rm bottom}\bigr)\ne0.
\]
Such a choice exists: use $z=0$ if the constant term is nonzero;
otherwise use $z=e_j$ for a nonzero coordinate of $cW$.

Set $g=Wz+Eb_{\rm top}-b_{\rm bottom}$ and
$y=x+Uz-b_{\rm top}$. The top inequalities require
\[
 y_i\in\bigl((Uz-b_{\rm top})_i+D\bigr)\setminus\{0\},
\]
a list with at least two elements. The bottom inequalities are
$Ey+g\in(\F^\times)^m$. Their cycle gain is nonzero by construction.
Color the unique cycle by Lemma~\ref{lem:cycle}, then color every
attached tree greedily away from the cycle. Each new vertex avoids
at most one value. Recovering $x=y-Uz+b_{\rm top}$ gives $x\in D^m$
and all the required inequalities.
\end{proof}

Theorem~\ref{thm:unicyclic} applies in arbitrarily large even dimensions.
It is a restricted affirmative class, not a proof of general DeVos
choosability. Its proof does not rely on numerical optimization.
The displayed block form is a structural hypothesis: general row
operations do not preserve coordinatewise nonzero constraints, and
general column operations do not preserve $D^{2m}$. It cannot be imposed
on an arbitrary instance by Gaussian elimination.

\paragraph{Finite searches and their limits.}
Additional heuristic searches up to dimension ten found no counterexample.
An integer-programming solver reported infeasibility in one dimension-four
encoding; this output was not converted into an independently checkable
exact infeasibility certificate. It is not used as a theorem here.
Enumeration of homogeneous forced-hyperplane constructions in dimensions
three and four found no translate-covering failure. Those finite results
also do not close the all-dimension problem. Theorems~\ref{thm:network}
and~\ref{thm:unicyclic} have the independent analytic proofs just given.

\section{Krzy\.{z}: an entropy-extension obstruction}

The coefficient target is
\begin{equation}\label{eq:krzyz}
 |a_n|\le\frac2e\quad(n\ge1),\qquad
 f(z)=\sum_{n\ge0}a_nz^n,\quad 0<|f(z)|\le1\ (z\in\D).
\end{equation}
The latest versions inspected are \cite{Wikstrom,MacDonald,LeiZhang}.
Wikstr\"om's November revision supplies the sixth and seventh coefficient
cases. MacDonald proves restrictions on global extremizers, including
$|f(0)|\le1/e$ and its equality characterization. These are prior
results and are not counted as progress from this report.

In the atomic extremal representation
\[
 f(z)=\exp\left(-\sum_{j=1}^{N}\lambda_j
                  \frac{1+\alpha_jz}{1-\alpha_jz}\right),
 \qquad |\alpha_j|=1,\quad \lambda_j>0,\quad N\le n,
\]
write $t=\sum_j\lambda_j$. The cited work gives
$\lambda_j\ge1/n$ and $t\ge n/N$ under its extremal assumptions.
Neither inequality by itself forces $N=n$. The entropy theorem in
\cite{LeiZhang} applies when every polynomial zero lies on the unit
circle; it does not remove the remaining full-degree issue.

The next exact calculation prevents a particular unjustified extension
of that entropy argument. It is \emph{not} a counterexample to
\eqref{eq:krzyz}.

\begin{proposition}\label{prop:entropy}
Let $dm$ be normalized arc measure on $\T$, and put
\[
 q(z)=(1-z)(1+z/2),\qquad p(z)=q(z)/\sqrt{3/2}.
\]
Then $\int_{\T}|p|^2\,dm=1$, its zeros are $1,-2$, and
\begin{equation}\label{eq:entropy}
 \int_{\T}|p|^2\log|p|^2\,dm
   =\frac7{12}-\log\frac32
   <1-\log2.
\end{equation}
\end{proposition}
\begin{proof}
The coefficients of $q=1-z/2-z^2/2$ give
$\int|q|^2\,dm=3/2$. Its squared modulus has Fourier coefficients
$c_1=-1/4$ and $c_2=-1/2$. Since
\[
 \log q(z)=-\frac z2-\frac{5z^2}{8}+O(z^3),
\]
Fourier pairing gives
\[
 \int_{\T}|q|^2\log|q|^2\,dm
 =2\left[\left(-\frac14\right)\left(-\frac12\right)
          +\left(-\frac12\right)\left(-\frac58\right)\right]
 =\frac78.
\]
This pairing can be justified first with radius $r<1$ in the logarithms
and then by their $L^1(\T)$ boundary convergence; the only boundary
singularity is logarithmic at $z=1$. Normalization gives
\eqref{eq:entropy}. Finally, the difference from $1-\log2$ is
$-5/12+\log(4/3)<-5/12+1/3<0$.
\end{proof}

Thus an entropy lower bound valid for polynomials whose zeros all lie
on $\T$ cannot simply be extended to polynomials with a boundary zero
and an exterior zero. The example has not been shown to satisfy the
extremal conditions of the Krzy\.{z} problem. A valid global extremal
argument remains necessary.

\section{What remains open in this work}

For Marcus--de Oliveira, neither a proof of \eqref{eq:moc} nor exact
separating data \eqref{eq:separation} have been produced. For DeVos,
the two structured classes above do not exhaust invertible matrices,
and no covering counterexample has been found. For Krzy\.{z}, no
all-index coefficient bound or global full-degree theorem has been
proved here.

The exact results in Sections 3--4 are construction barriers and
restricted affirmative statements. They do not meet the requirement
of a new unconditional resolution of an unchanged, previously open,
named problem. A bounded literature review also cannot certify
historical priority. These limits are part of the report's conclusion,
not assumptions to be silently discarded in a later manuscript.

\begin{thebibliography}{9}
\bibitem{MOCsurvey}
N.~Bebiano and J.~P.~da Provid\^encia,
\emph{Revisiting the Marcus--de Oliveira Conjecture},
Mathematics \textbf{13} (2025), 711.
\url{https://doi.org/10.3390/math13050711}.
\bibitem{Holden}
S.~Holden, \emph{MI-05: A fourteen-term identity and a sharp orthogonal
obstruction bound}, research report, 13 September 2026.
\url{https://github.com/ajt60gaibb/OpenProblemsInNLA/blob/eb7c9d64f91757a82fbe9167e80ab9c1b7f36c22/references/holden-mi05-2026-09-13/report.pdf}.
\bibitem{DeVos}
M.~DeVos, author-posted formulation of the choosability conjecture,
\emph{A nowhere-zero point in a linear mapping}, Open Problem Garden,
posted 8 March 2007.
\url{https://www.openproblemgarden.org/op/a_nowhere_zero_point_in_a_linear_mapping}.
\bibitem{Wikstrom}
F.~Wikstr\"om, \emph{A computer-assisted proof of Krzy\.{z}'s conjecture
for the sixth coefficient}, arXiv:2610.15206v2, 4 November 2026.
\url{https://arxiv.org/abs/2610.15206v2}.
\bibitem{MacDonald}
S.~F.~MacDonald, \emph{A convolution inequality and the Krzy\.{z}
conjecture}, arXiv:2611.05363v2, 5 November 2026.
\url{https://arxiv.org/abs/2611.05363v2}.
\bibitem{LeiZhang}
J.~Lei and T.~Zhang, \emph{Proof of the Agler--McCarthy entropy
conjecture}, arXiv:2605.03949v2, 6 May 2026.
\url{https://arxiv.org/abs/2605.03949v2}.
\end{thebibliography}

\clearpage
\appendix
\section{Standalone exact checks of the finite-field lemma}

The following Python 3 program uses only the standard library. It checks
the determinant, the six survivors, the sixteen-point plane section,
every translate in Lemma~\ref{lem:translate}, every affine offset in
Lemma~\ref{lem:local}, and the rational Fourier calculation. It does
not purport to verify the original conjectures. The arbitrary-dimensional
arguments are the analytic proofs above.

\begin{lstlisting}
from itertools import product
from fractions import Fraction as Q

p = 5
D = (0, 1, 4)
A = ((2, 0, 2), (1, 4, 0), (3, 3, 2))
cube = list(product(D, repeat=3))
space = list(product(range(p), repeat=3))

def require(condition, message):
    if not condition:
        raise ArithmeticError(message)

def mv(M, x):
    return tuple(sum(a*b for a, b in zip(row, x)) % p
                 for row in M)

def L(x):
    return (x[1] + x[2]) % p

def nonzero(x):
    return all(a != 0 for a in x)

def inverse(M):
    n = len(M)
    T = [[a % p for a in M[i]] +
         [int(i == j) for j in range(n)] for i in range(n)]
    for j in range(n):
        k = next(i for i in range(j, n) if T[i][j])
        T[j], T[k] = T[k], T[j]
        s = pow(T[j][j], -1, p)
        T[j] = [(s*a) % p for a in T[j]]
        for i in range(n):
            if i != j:
                s = T[i][j]
                T[i] = [(a-s*b) % p for a, b in zip(T[i], T[j])]
    return [row[n:] for row in T]

det = (A[0][0]*(A[1][1]*A[2][2]-A[1][2]*A[2][1])
       - A[0][1]*(A[1][0]*A[2][2]-A[1][2]*A[2][0])
       + A[0][2]*(A[1][0]*A[2][1]-A[1][1]*A[2][0]))
require(det == -2, "integer determinant")
Ai = inverse(A)
require(tuple((Ai[1][j]+Ai[2][j]) % p for j in range(3))
        == (3, 4, 0), "L times inverse")

survivors = {x for x in cube if nonzero(mv(A, x))}
expected = {(a, b, (-b) % p) for a in D for b in D if a != b}
require(survivors == expected and len(survivors) == 6,
        "six survivors")
K = {x for x in space if L(x) == 0 and nonzero(mv(A, x))}
require(len(K) == 16, "plane section cardinality")

for h in space:
    translated = {tuple((h[i]+x[i]) % p for i in range(3))
                  for x in cube}
    require(bool(K & translated), "translate cover")

for a in space:
    target = (-L(mv(Ai, a))) % p
    candidates = [x for x in cube if L(x) == target and
                  nonzero(tuple((v+a[i]) % p
                                for i, v in enumerate(mv(A, x))))]
    require(bool(candidates), "affine local realization")

coeff = (Q(1), Q(-1, 2), Q(-1, 2))
norm2 = sum(a*a for a in coeff)
c1 = coeff[1]*coeff[0] + coeff[2]*coeff[1]
c2 = coeff[2]*coeff[0]
raw_entropy = 2*(c1*Q(-1, 2) + c2*Q(-5, 8))
require(norm2 == Q(3, 2), "squared L2 norm")
require(raw_entropy == Q(7, 8), "raw entropy integral")
require(raw_entropy/norm2 == Q(7, 12), "normalized rational part")
print("PASS: exact local lemmas and rational entropy calculation.")
print("No original named conjecture is resolved by these checks.")
\end{lstlisting}
\end{document}
